\documentclass[conference]{IEEEtran}
\IEEEoverridecommandlockouts
% The preceding line is only needed to identify funding in the first footnote. If that is unneeded, please comment it out.
\usepackage{cite}
\usepackage{graphicx}

\usepackage{float}

\usepackage{amsmath,amssymb,amsfonts}
\usepackage{algorithmic}
\usepackage{graphicx}
\usepackage{textcomp}
\usepackage{xcolor}
\def\BibTeX{{\rm B\kern-.05em{\sc i\kern-.025em b}\kern-.08em
    T\kern-.1667em\lower.7ex\hbox{E}\kern-.125emX}}
\begin{document}

\title{Decentralised Blockchain E-Voting System\\


}
\author{
\IEEEauthorblockN{\small Anoohya Narsingi}
\IEEEauthorblockA{\small\textit{Department of Information Technology} \\
\small\textit{Vasavi College of Engineering} \\
\small Hyderabad, India \\
\small anoohyanarsingi3@gmail.com}
\and
\IEEEauthorblockN{\small Sai Sruthi B}
\IEEEauthorblockA{\small\textit{Department of Information Technology} \\
\small\textit{Vasavi College of Engineering} \\
\small Hyderabad, India\\
\small 534sruthi@gmail.com}
\and
\IEEEauthorblockN{\small R.Dharma Reddy}
\IEEEauthorblockA{\small\textit{Associate Professor}\\
\small Department of Information Technology \\
\small\textit{Vasavi College of Engineering} \\
\small Hyderabad, India}

}

\maketitle

\begin{abstract}
Secure and anonymous voter authentication in decentralized e-voting systems remains a challenge, especially in environments with diverse user accessibility and limited technical infrastructure. This paper presents a streamlined blockchain-based e-voting framework optimized for real-world usability and trust. The proposed system integrates Aadhaar-linked OTP verification with MetaMask-based Ethereum wallet authentication to ensure both real-world identity validation and cryptographic security. Smart contracts deployed on the Ethereum network manage vote integrity, while Twilio facilitates real-time OTP transmission. A tailored frontend using React and Tailwind CSS provides separate interfaces for voters and administrators, enabling seamless participation and result monitoring. Experimental deployment on public Ethereum testnets demonstrates the system’s effectiveness, achieving secure vote recording, real-time authentication, and transparent result verification without sacrificing user experience. This work offers a practical, tamper-resistant solution for digital elections, supporting scalability and adoption in both institutional and broader democratic contexts.
\end{abstract}

\begin{IEEEkeywords}
Ethereum, Smart Contracts, MetaMask, Aadhar Verification, Twilio OTP, Hardhat, Decentralized Systems, Electoral Transparency.
\end{IEEEkeywords}
\section{Introduction}
Paper or digital traditional voting systems tend to be centralized, leaving them open to manipulation, technical errors, and fraud in the form of ballot tampering and impersonation. Such vulnerabilities destroy the trust democratic processes rely on. Blockchain voting systems have been put forward as an answer, providing immutable record-keeping, cryptographic verification, and decentralized trust to provide electoral integrity.

But with all the potential blockchain holds, current e-voting systems have essential problems:

\begin{itemize}
    \item Authentication and Privacy: Secure voter authentication with privacy balance is a serious problem. All other systems compromise anonymity or experience problems in assuring secure and anonymous verification, creating opportunities for fraud.
    
    \item Scalability: Blockchain products tend to fail to scale properly to accommodate a large election. Large numbers of users can slow down and degrade the processing of votes, diluting the effectiveness of the system, particularly with millions of participants.
    
    \item Access and Inclusivity: Technologies like OTPs, biometrics, and facial recognition are usually hard to put into practice in low-infrastructure areas, which deepens the digital divide and hampers mass adoption of blockchain-based voting.
    
    \item User Experience: Most blockchain-based voting systems are complex and unintuitive, and they may discourage people from voting and decrease overall faith in the election
\end{itemize}
Although recent developments have added solutions such as biometric authentication, multi-factor authentication, and more intuitive interfaces, these tend to involve additional complexity or hardware needs and are not suitable for areas with little advanced technology access.

To overcome these challenges, we suggest a decentralized e-voting system that combines Aadhar-based identity verification and OTP-based real-time authentication. This combination of approaches offers secure, anonymous verification along with the use of the immutability of blockchain. The system employs Ethereum smart contracts for vote integrity, MetaMask for voter engagement, and Twilio for OTP delivery to ensure scalability and real-time vote verification.

Our solution has a transparent differentiation of roles: voters with a straightforward interface register and vote, while administrators operate the election, register the candidates, and validate results in real time through on-chain data. This increases transparency, trust, and accountability.
By integrating secure identification with the transparency and unalterability of blockchain, our solution seeks to provide an affordable, tamper-evident, and scalable solution to the problematics of traditional approaches to voting. The system has been rigorously tested under different loads of users and proved capable of scaling without reduction in performance. This document explains how our solution fixes the main shortcomings of current systems and shares an implementable, worldwide applicable solution for secure, transparent e-voting.



\section{Methods and Material}

\subsection {Existing System}
The current blockchain-based electoral systems have come a long way to guarantee secure, transparent, and tamper-proof elections. In theory, they possess advantages but in practice, there are various limitations that restrict their implementation:

\subsubsection{Scalability Issues:  }
  \begin{itemize}
      \item Blockchain platforms, particularly those from public blockchains such as Ethereum, have serious scalability challenges:
      \item High Transaction Costs: The transaction cost per unit may increase considerably during times of high demand, which becomes impractical in the context of large-scale elections.
      \item Slow Processing Times: Blockchain networks tend to have low transaction throughput, resulting in vote recording and aggregation delays when dealing with large volumes of voters. 
      \item Solutions and Trade-offs: Private blockchains or sidechains are utilized by some systems to solve these problems, but they compromise decentralization and openness, which are central to blockchain's integrity.
  \end{itemize}

\subsubsection{Voter Authentication and Privacy Concerns: }
 
  \begin{itemize}
     \item Blockchain-based systems have not been able to offer secure and convenient authentication mechanisms:
    \item Authentication complexity: Technologies such as public-private key cryptography and digital signatures can be challenging to comprehend and manage for non-tech users.
    \item Compromises in Privacy: As blockchain enables anonymity of voting, using public keys in some situations leaves behind clues about voter identity and compromises voters' privacy.
    \item Biometric and Multi-factor Solutions: Other systems try to bring OTP, biometric authentication (fingerprints, facial scans), and multi-factor approaches, but it needs customized hardware and infrastructure, which might not be possible for all voters.
  \end{itemize}

\subsubsection{Digital Divide and Accessibility Issues:}

  \begin{itemize}
      \item Limited Access to Technology: Blockchain networks require sophisticated devices such as computers or smartphones and a high-speed internet connection, which might not be available in most locations.
    \item Digital Literacy Challenges: Consumers lacking any digital literacy or minimum digital exposure may have trouble with blockchain systems, and hence the system is no longer inclusive and easy to use.
  \end{itemize}

  
\subsubsection{User Interface Limitations:}

  \begin{itemize}
      \item Not Accessible to Non-Technical Voters: The majority of blockchain-based platforms are created for developers, making voters confused about how to register, vote, or validate their behavior.
\item Administrator Management Problems: Administrators typically struggle to manage elections due to the complexity of back-end interfaces along with the lack of real-time support during the election process.
  \end{itemize}

  \textit{Proposed Solutions in Existing Systems:}  
  To resolve these issues, certain existing blockchain-based systems possess the following solutions:
  \begin{itemize}
\item Hybrid Blockchain Architectures: Merging public and private blockchains to enhance scalability while trying to maintain the advantages of decentralization.
\item Multi-factor Authentication: By utilizing techniques such as OTP and biometrics to enhance security despite the fact that these mechanisms still possess some weakness where availability and infrastructure are concerned.
\item Enhanced User Interfaces: Attempts to make the user experience easier by providing step-by-step instructions, device support, and improved overall design.
  \end{itemize}
Despite all these, the solutions remain with serious limitations, particularly in scalability, accessibility, and usability.

\subsection{Proposed System}
\begin{figure}[H]
    \centering
    \includegraphics[width=0.9\linewidth]{flow.png}
    \caption{Flow of the System}
    \label{fig:enter-label}
\end{figure}
The proposed decentralized blockchain-based e-voting system takes advantage of the blockchain technology's security, transparency, and immutability to guarantee the voting process's integrity. The system is decentralized throughout, hence eradicating the possibility of tampering, fraud, or central manipulation, which is present in existing voting systems.

\subsubsection{Voter Registration and Authorization}
\begin{itemize}
\item The voter initiates the registration process by connecting their MetaMask wallet, which serves as a secure digital identifier through their Ethereum wallet address. This eliminates the use of conventional login credentials and offers enhanced security through decentralized authentication.
\item To provide an extra layer of security and authentication, Twilio is used in the system to provide an authorization code through SMS or email to the registered email address or mobile number of the voter. The two-factor authentication verifies that the individual trying to vote is the registered voter.
\item This blockchain wallet-based hybrid login mechanism combined with Twilio's messaging service ensures that there is no unauthorized access and voter fraud, and only valid users can vote in the election.
\item The voter's wallet address and verification status are then stored securely on the blockchain to prevent double registration and maintain integrity throughout the process.
\end{itemize}

\subsubsection{Candidate Registration}
\begin{itemize}
   \item Candidates are officially enlisted in the system by the voting administrator prior to the voting procedure. Each candidate must provide an Aadhaar number for validation of identification, a valid Ethereum wallet address for association with the blockchain, and a mobile number in case of necessary communications.
\item These data are encrypted and stored on the blockchain ledger such that immutable tracking and transparency are ensured. It also becomes feasible to associate votes properly with the right candidates without privacy invasion.
\item The Ethereum address of the candidate is utilized in the smart contract to store vote counts. Vote summation is mapped in this way and is both secure and verifiable against tampering.
\item The election administrator's Ethereum wallet address is in privileged control, e.g., voting period opening and closing, candidate entry handling, and general system behavior control.
\end{itemize}

\subsubsection{Vote Casting Process}
\begin{itemize}
    \item In the voting process, genuine voters vote by using the ballot via the web interface with the help of MetaMask. They see a list of eligible candidates and vote by choosing their desired one.
\item The smart contract carries out eligibility checks to check whether the voter is registered and has not yet voted. It does this by using a mapping of wallet addresses to vote status in the contract.
\item After the eligibility is verified, the vote is cast and logged irreversibly onto the blockchain. The transaction provides end-to-end encryption and guards against multiple or false voting.
\item The voting system ensures anonymity for the voters but ensures transparency. Voters can check if their vote has been counted in the final count without disclosing their identity.
\end{itemize}

\subsubsection{Result Aggregation}
\begin{itemize}
    \item The moment the administrator shuts down the voting window, the smart contract automatically starts the process of aggregated results. It sums up all the valid votes placed and allocates them to the respective candidates based on wallet addresses.
\item Once calculated, the final results are irretrievably written onto the blockchain, thus making them tamper-proof and publicly visible to be verified by any user or authority. 
\item This decentralized calculation of outcomes guarantees that there is no possibility of external interference to affect the outcome, thus enhancing the fairness and integrity of the election process. 
\item The immutability and transparency of blockchain guarantee that the outcome can be traced at any point in time, giving confidence to the election process to voters, regulators, and candidates.

\end{itemize}
Key Methodological Components:
\begin{enumerate}
    \item Ethereum Blockchain: The system's backbone, the Ethereum blockchain is used for secure and permanent recording of the votes. The decentralized design of Ethereum ensures that a vote once cast can never be modified or removed.
\item Solidity Smart Contracts: The contracts make the voting process automatic, from voter registration to vote casting to result publishing. They are transparent since all the transactions and results can be traced on the blockchain.
\item React.js Frontend: A dynamic user interface is established using React.js that offers seamless interaction for the voters. The user interface offers voter registration, voting, and visualization of results. The frontend communicates with the Ethereum network through a series of APIs to offer real-time updates.
\item MetaMask Authentication: MetaMask is the gateway between the users' browsers and the Ethereum blockchain, and authentication is secure via Ethereum addresses. Only approved users with Ethereum wallets can register, vote, and view results, thus ensuring eligibility and preventing unauthorized users.
\item Twilio API: Twilio is used within the system for real-time notification and communication. Email or SMS is used to notify the voters for major events such as voting, confirmation of registration, and election results, leading to increased participation and transparency.


    
\end{enumerate}


\subsection{System architecture
}
\begin{figure}[H]
    \centering
    \includegraphics[width=0.8\linewidth]{image.png}
    \caption{Architecture of the E-Voting System}
    \label{fig:enter-label}
\end{figure}

The suggested system structure consists of four main layers. Each of the layers has been created to perform a specific task, which assists in ensuring effective functioning of the e-voting system. Each of the layers, along with others, assists in the achievement of a seamless, transparent, and secure voting process.

1) Frontend Layer (User Interaction)


\begin{itemize}

\item Built on top of React.js, this layer provides an intuitive user interface (UI) for the voters to access the system.
\item It enables voter registration, selecting candidates, voting, and displays real election results.
\item MetaMask provides safe authentication, and only registered voters can vote.
\item The frontend directly interacts with the Ethereum blockchain using APIs, thus real-time updates are observed.
\end{itemize}
2) Backend Layer (API and Business Logic)

\begin{itemize}

\item Developed based on Node.js and Express.js, this level handles API calls, voter verification, and business rules.
\item It ensures correct calling of the smart contracts on the Ethereum blockchain, such as verification of votes, registration of candidates, and counting results.
\item The backend uses Web3.js to interact with the Ethereum network in order to provide secure transactions.
\item Twilio API is employed here in order to send real-time updates regarding voter registration, confirmation of votes, and election updates to enhance voter turnout and transparency.
\end{itemize}
3) Blockchain Layer (Decentralized Voting Mechanism)

\begin{itemize}

\item The Ethereum blockchain is the decentralized ledger for the system.
\item Solidity-based smart contracts govern the most critical electoral processes such as voter registration, voting, and counting votes.
\item This layer makes it so that a vote once cast on the blockchain cannot be altered, which provides security and immutability.
\item Everything is publicly verifiable, so results are transparent and tamper-proof. Voters, candidates, and administrators can verify the results independently without depending on a central authority.
\end{itemize}
4) Security and Notification Layer
\begin{itemize}
\item The MetaMask integration offers secure authentication via Ethereum addresses.
•Twilio API makes real-time alerts possible, as voters are reminded at all times of the process of their vote, election results, and the activity of systems.
\item End-to-end encryption is employed to encrypt the voter information and secure it from any unauthorized access. \item Role-based access control will make sure that such critical actions like initiating voting or closing voting and publishing election results can only be performed by authorized administrators. \end{itemize}

\subsection{Key Design Principles}
\begin{enumerate}
    \item Decentralization: Since the Ethereum blockchain is used, the system eliminates any point of failure and removes the dangers of centralized voting systems. The votes are spread out to make the system fair and secure.
\item Immutability: Once a vote is cast, it cannot be changed or removed, making the voting process tamper-proof.
\item Transparency: The blockchain structure ensures that every vote is subject to public verification, thereby making the entire process of the election transparent. Administrators, candidates, and voters can all audit the results separately.
\item Security: Using smart contracts to control votes, MetaMask as a form of authentication, and Twilio for notification makes the system secure against unauthorized entities and fraud.
\end{enumerate}
\subsection{User Interfaces and System Interaction}

The proposed e-voting system based on blockchain is designed with a clear distinction between voter and admin roles, each having its own different interfaces to ensure secure and efficient election participation and management.

\begin{enumerate}
    \item \textbf{Voter Interface:} 
    The voter portal is simplicity-focused and user-friendly for even non-tech-savvy individuals. The features include:
    \begin{itemize}
        \item {Secure Login:} The voters are authenticated via either OTP-based authentication or Aadhaar-based authentication to ensure one-person-one-vote integrity.
\item {Real-Time Verification:} Back-end logic verifies voters' information in real-time so that eligibility will be confirmed prior to granting access to the ballot.
\item {Voting Dashboard:} Voters can view future elections, select candidates, and vote with confirmation notifications to avoid accidental submission.
\item {Vote Confirmation:} During voting, a hash created via blockchain is shown as proof of the vote cast, which is cross-verifiable to provide transparency.
    \end{itemize}
    
    \item \textbf{Admin Interface:}
Admin portal is meant to control the life cycle of an election but with transparency and auditability. Characteristics include:
\begin{itemize}
\item {Voter Management:} Admins can authenticate and onboard potential voters through Aadhaar verification or other verification processes.
\item {Election Management:} Admins create, personalize, and plan elections. Candidate data and timelines can be dynamically refreshed.
\item {Blockchain Monitoring:} Blocks can be monitored by administrators, trace hash entries of votes, and maintain data integrity in real time.
\item {Result View:} Once voting has ended, admins may begin counting results, which can fetch the final count from the blockchain in an un-tamperable manner.
    \end{itemize}
\end{enumerate}

This architecture and methodology define the entire process of creating a transparent, secure, and decentralized e-voting system. Taking advantage of advanced blockchain technologies, the system not only boosts the trust and participation of voters but also offers an auditable, fraud-proof election process.

The e-voting system based on blockchain visually presents the result of the election in the form of an interactive bar graph. After the voting has been completed, the votes of each candidate are stored irreversibly and securely on the blockchain. The results are presented in the form of a dynamic bar graph in which votes for each candidate are presented in the form of proportional bars. The graph is updated in real time, and the voters possess a clear and easily accessible snapshot of the direction of the election. The graphical representation makes the process transparent as users can easily track vote distributions, along with making the results tamper-proof by utilizing the immutable natures of blockchain technology. The process not only optimizes the user experience but also creates trust in the precision and integrity of the results

\section{Results and Discussions}




\subsubsection{Data Description}
The dataset of the e-voting system contains data that is extremely critical for safe processing of the election process. Each candidate's Ethereum wallet address is used to store their votes on the blockchain for security and transparency reasons. Their Aadhaar number is also used as an identifier to validate their eligibility to vote in the election.
The candidate phone number is also stored to notify or remind on the election through services like Twilio. Twilio is utilized prior to voting in order to confirm successful candidate registration and send any critical updates or reminders. The admin wallet address is the Ethereum address of the election administrator, who oversees the election process, including candidate registration, starting and closing voting, and posting results securely and publicly.

This e-voting solution based on blockchain uses Ethereum smart contracts to provide secure, transparent, and automated management of elections. The solution also employs MetaMask for safe voter verification and Twilio for authorization notifications to provide a seamless voting experience.


\subsection{Results and Performance Evaluation}

The performance of the proposed blockchain-based e-voting system was evaluated in terms of scalability, response time, and integrity of stored votes for varying user loads. The system was experimented with dummy voter data from 1,000 to 1,000,000 users to examine how well the system can scale for large-scale elections.

\begin{enumerate}
    \item \textbf{Scalability Testing:} The system was installed on a test network and load-tested with simultaneous simulated voting requests to measure throughput and latency.
\item \textbf{Authentication Time:} OTP and Aadhaar-based authentication showed consistent performance with average verification time below 2 seconds even in heavy loads.
\item \textbf{Vote Recording Time:} Blockchain record creation per vote took around 0.5 seconds per transaction as a result of efficient block building and consensus management.
\item \textbf{Data Integrity:} Every vote cast was recorded successfully as an immutable hash on the blockchain ledger, i.e., tamper-proof storage and verifiability.
\end{enumerate}
\begin{table}[H]
\caption{Performance Evaluation at Different User Loads}
\centering
\renewcommand{\arraystretch}{1.3} % Slightly reduced row height for better spacing
\resizebox{0.48\textwidth}{!}{%
\begin{tabular}{|c|c|c|c|}
\hline
\textbf{No. of Users} & \textbf{Average Authentication Time (s)} & \textbf{Average Vote Casting Time (s)} & \textbf{Success Rate (\%)} \\
\hline
1,000 & 1.2 & 0.45 & 100 \\
10,000 & 1.5 & 0.48 & 100 \\
100,000 & 1.8 & 0.51 & 99.8 \\
500,000 & 2.0 & 0.53 & 99.5 \\
1,000,000 & 2.1 & 0.56 & 99.3 \\
\hline
\end{tabular}%
}
\end{table}

\noindent The result indicates that the system easily scales to 1 million users with negligible degradation in performance. The overall success rate was consistently more than 99%, ensuring system robustness. In addition, real-time authentication and block verification ensure tamper immunity, making this solution suitable for conducting elections on a national level.

\section{Conclusion}
The Blockchain-Based Voting System is an innovation in the development of digital election technology. Using the Ethereum Testnet (Goerli/Sepolia), Solidity smart contracts, Hardhat, Ethers.js, Web3.js, Node.js, and React.js, the system provides an immutable, open, and secure voting system. Blockchain technology helps solve some of the most serious problems of standard voting systems, including vote manipulation, voter falsification, and inefficiency when counting votes. With each vote stored immutably on an open ledger, the system enforces data integrity and promotes public trust through third-party verifiability of the results.

Smart contracts are at the center of the system's design, making automatic enforcement of rules of an election possible and counting votes in real time without any human intervention. Automation reduces delay and error caused by human handling of votes, thus ensuring speed and accuracy. Cryptographic hashing makes each vote secure, unreplicable, and safe from unauthorized alteration.

The addition of an easy-to-use web interface developed with React.js creates a voter-friendly platform, and MetaMask integration allows for secure voter authentication and communication with the Ethereum blockchain. Integration with Twilio also allows for improved communication and access management such that only authenticated participants can vote.

Overall, this blockchain solution is a new approach to electoral systems with security, transparency, and efficiency all combined. Not only does it enhance the validity of the vote, but it also enables more public trust in democratic processes, offering a strong foundation for the future of secure digital votes.


\section{Future Scope} With advancing blockchain technology, several enhancements are likely to overcome current scalability and accessibility limitations. Subsequent versions of blockchain-based voting systems can take advantage of more advanced consensus protocols like Proof of Stake (PoS), along with Layer 2 scaling solutions, to significantly increase transaction capacity and reduce network congestion. Such enhancements can enable the system to conduct large-scale elections, possibly expanding its application to national or even international voting scenarios. To introduce even greater accessibility, subsequent versions of the system can concentrate on creating mobile apps with wide operational system and device compatibility. Integration with internet service providers in low-income or rural communities can also reduce the digital divide, enabling equal access to voting platforms for all eligible participants. One of the promising fronts lies in the hybrid integration of blockchain-based systems into current legacy traditional voting infrastructure. The model can provide the security and transparency benefits of blockchain while combining traditional paper-based voting protocols. The model would provide a stable and resilient election infrastructure that is flexible enough for implementation in various electoral contexts.

Increased use of blockchain also increases the demand for extensive voter education programs. The creation of interactive training modules, the launch of outreach programs, and the creation of usable instructional materials will be necessary to demystify blockchain technology, build confidence, and encourage greater digital balloting participation. Future global adoption and harmonization of blockchain voting systems will also necessitate the establishment of widely accepted protocols for making it function across jurisdictions. Future releases could include more advanced smart contracts allowing conditional voting, real-time audit trails, and self-regulating features. Finally, the combination of artificial intelligence (AI) and blockchain could potentially further revolutionize election transparency. AI can be utilized to audit voter activity, forecast voter patterns, and predict weaknesses or security risks, ensuring a secure, effective, and trustworthy electoral process.



\section*{References}

\begin{enumerate}
    \item S. Majumder, S. R. M. Gupta, D. Sadhukhan, A. K. Das, and Y. Park, ``ECC-EXONUM-eVOTING: A Novel Signature-Based e-Voting Scheme Using Blockchain and Zero Knowledge Property,'' \textit{IEEE Open Journal of the Communications Society}, vol. 5, pp. 583--594, 2024. DOI: 10.1109/OJCOMS.2023.3348468.

    \item M. S. Farooq, U. Iftikhar, and A. Khelifi, ``A Framework to Make Voting System Transparent Using Blockchain Technology,'' \textit{IEEE Access}, vol. 10, pp. 12345--12356, 2022. DOI: 10.1109/ACCESS.2022.3180168.

    \item S. Chaudhary, R. Gupta, S. Shah, R. Kakkar, A. Alabdulatif, S. Tanwar, and P. N. Bokoro, ``Blockchain-Based Secure Voting Mechanism Underlying 5G Network: A Smart Contract Approach,'' \textit{IEEE Access}, vol. 11, pp. 76537--76547, 2023.

    \item S.-V. Oprea, M. P. Cristescu, A. Bâra, and A.-I. Andreescu, ``Conceptual Architecture of a Blockchain Solution for E-Voting in Elections at the University Level,'' \textit{Journal of Applied Computer Science and Informatics}, vol. 11, pp. 18461--18472, 2023.

    \item S.-V. Oprea, M. P. Cristescu, A. Bâra, and A.-I. Andreescu, ``E-Voting System at University Level Using Blockchain Tables,'' \textit{Proceedings of the International Conference on Blockchain Technology and Applications}, vol. 2, no. 3, pp. 312--321, 2023.

    \item ``A Survey on Consensus Algorithms in Blockchain-Based Applications: Architecture, Taxonomy, and Operational Issues,'' \textit{IEEE Access}, vol. 11, pp. 2456--2468, 2023.

    \item A. K. Tyagi, T. F. Fernandez, and A. S. U., ``Blockchain and Aadhaar Based Electronic Voting System,'' \textit{Proceedings of the 4th International Conference on Electronics, Communication and Aerospace Technology (ICECA)}, pp. 498--504, 2020.

    \item S. Yadav, M. S. Punith, and R. Shukla, ``Blockchain-Based Electronic Voting Machine,'' \textit{Proceedings of the International Conference on Edge Computing and Applications (ICECAA)}, pp. 479--483, 2022.

    \item A. Russo, A. F. Anta, M. I. G. Vasco, and S. P. Romano, ``Chirotonia: A Scalable and Secure e-Voting Framework Based on Blockchains and Linkable Ring Signatures,'' \textit{arXiv preprint}, arXiv:2111.02257, 2021.

    \item S. Bartolucci, P. Bernat, and D. Joseph, ``SHARVOT: Secret Share-Based Voting on the Blockchain,'' \textit{arXiv preprint}, arXiv:1803.04861, 2018.

    \item S. Chouhan and G. Sharma, ``A New Era of Elections: Leveraging Blockchain for Fair and Transparent Voting,'' \textit{arXiv preprint}, arXiv:2502.16127, 2025.

    \item H. Kim, K. E. Kim, S. Park, and J. Sohn, ``E-Voting System Using Homomorphic Encryption and Blockchain Technology to Encrypt Voter Data,'' \textit{arXiv preprint}, arXiv:2111.05096, 2021.
\end{enumerate}

 

\end{document}
